\chapter{Broken Toe}

\section*{Definition and Overview}
A broken toe, or toe fracture, is a crack or break in one of the small bones that form the toes. Toe fractures are among the most common fractures of the foot and are typically caused by dropping a heavy object on the foot or stubbing the toe against a hard surface. The great (big) toe contains two bones (phalanges), while each of the other toes has three. The majority of broken toes are minor injuries that heal well with simple home care, but some require specific treatment, particularly fractures of the big toe, displaced fractures, breaks that involve a joint, or open fractures where the skin is broken. Distinguishing a fracture from a bad bruise or sprain matters because it guides treatment and follow-up.

\section*{Epidemiology}
Toe fractures are extremely common in adults and slightly more frequent in men, reflecting higher rates of workplace injuries. Typical causes include stubbing the toe on furniture, dropping heavy objects on the foot, and kicking hard surfaces accidentally. They are also common in sport, particularly football, running, and activities involving rapid stops and turns. In older adults, toe fractures often follow trips and falls and may be related to reduced bone density. Many toe fractures are never brought to medical attention, so the true incidence is unknown, but they account for a substantial share of all foot injuries seen in emergency departments.

\section*{Aetiology and Pathophysiology}
\begin{itemize}
    \item \textbf{Direct impact:} dropping a heavy object on the toe or a crush injury is the most common mechanism
    \item \textbf{Axial loading:} stubbing the toe or kicking a hard object drives force along the length of the bone, causing a transverse or comminuted break
    \item \textbf{Twisting:} catching the toe and twisting the foot can fracture or dislocate the toe
    \item \textbf{Stress fracture:} repeated loading from running or prolonged walking can produce a hairline crack, most often in the second or third toe
    \item \textbf{Bone weakness:} osteoporosis or other metabolic bone disease increases fracture risk in older people
    \item \textbf{Pathophysiology:} the break damages bone, vessels, and soft tissues, causing bleeding and swelling; in displaced fractures the bone ends separate and the toe may appear crooked
\end{itemize}

\section*{Clinical Features}
\begin{itemize}
    \item Immediate pain in the toe, aggravated by walking, pressure, or weight-bearing
    \item Swelling of the toe and often of the forefoot
    \item Bruising that may spread widely, sometimes tracking under the toenail
    \item Localised tenderness over the broken bone
    \item A crooked, shortened, or angulated appearance when the fracture is displaced
    \item Difficulty weight-bearing or wearing a closed shoe
    \item Big toe fractures cause more pain and functional impairment than fractures of the smaller toes
    \item Numbness or tingling from nerve injury, or a pale, cold toe from impaired circulation (rare)
\end{itemize}

\section*{Differential Diagnosis}
\begin{itemize}
    \item \textbf{Sprain or ligament injury:} tearing of the ligaments around the toe joints, causing similar pain and swelling
    \item \textbf{Contusion (bruising):} soft-tissue injury without a fracture
    \item \textbf{Dislocation:} a toe joint displaced out of position, which looks markedly crooked
    \item \textbf{Subungual haematoma:} a painful blood collection beneath the toenail following injury
    \item \textbf{Metatarsal fracture:} a break in the long bones of the foot, causing similar pain and difficulty walking
    \item \textbf{Gout or infection:} a hot, red, swollen toe may be due to crystal arthritis or soft-tissue infection rather than trauma
    \item \textbf{Stress fracture:} a gradually developing ache in a toe or foot bone from overuse
\end{itemize}

\section*{When to Seek Medical Care}
See a doctor if the toe is very painful, the big toe is involved, the toe looks crooked, you cannot weight-bear, or pain does not settle within a few days of simple care. People with diabetes, poor circulation, or reduced sensation in the feet should have any foot injury checked, as healing is slower and infection more likely.

\textbf{Call 000 (or local emergency services) for:}
\begin{itemize}
    \item A toe pointing at an odd angle, or an open fracture where bone is visible through the skin
    \item A toe that becomes pale, blue, cold, or numb
    \item Severe pain not relieved by simple analgesia
    \item An open wound with heavy bleeding, or spreading redness with fever suggesting infection
    \item Numbness or weakness in the foot or toes after the injury
\end{itemize}

\section*{Diagnosis and Investigations}
\begin{itemize}
    \item \textbf{History:} the mechanism of injury, severity of pain, and ability to weight-bear
    \item \textbf{Examination:} inspection for swelling, bruising, and deformity; gentle palpation to localise the break; and assessment of the circulation and sensation in the toe
    \item \textbf{X-ray:} a plain radiograph confirms the fracture, reveals displacement or joint involvement, and excludes a dislocation
    \item \textbf{Whole-foot assessment:} checking for other fractures, since a force strong enough to break a toe can injure neighbouring bones
    \item \textbf{Risk-factor review:} questions about diabetes, peripheral vascular disease, smoking, and bone health
    \item \textbf{Repeat imaging:} may be arranged if healing is slow or symptoms persist
\end{itemize}

\section*{Management}
\begin{itemize}
    \item \textbf{Buddy strapping:} taping the injured toe to the adjacent toe with soft padding between to hold it straight and protect it; the standard treatment for fractures of the smaller toes
    \item \textbf{Firm footwear:} a stiff-soled, closed shoe or rigid-soled boot to reduce bending of the foot during walking
    \item \textbf{First aid:} rest, elevation, ice wrapped in a cloth applied for 15--20 minutes at a time, and avoidance of weight-bearing in the early phase
    \item \textbf{Pain relief:} simple analgesics such as paracetamol and ibuprofen
    \item \textbf{Big toe fractures:} may require a walking boot or short period of casting, as the big toe bears much of the body weight in gait
    \item \textbf{Displaced fractures:} a doctor may realign the bone (closed reduction); if this is not possible or fails, surgery with small pins or wires may be required
    \item \textbf{Subungual haematoma:} drainage of a painful blood blister beneath the nail may be offered
    \item \textbf{Follow-up:} review, and occasionally repeat X-ray, if pain is slow to settle
\end{itemize}

\section*{Complications}
\begin{itemize}
    \item Slow healing or non-union, more likely in smokers and those with poor circulation or diabetes
    \item Malunion, in which the bone heals crooked and leaves a permanent bump or slightly misaligned toe
    \item Post-traumatic arthritis in the affected joint, causing long-term stiffness and aching
    \item Osteomyelitis if an open wound becomes infected
    \item Nail injury, sometimes leaving permanent nail deformity
    \item Nerve damage with residual numbness
    \item Chronic pain or stiffness, particularly of the big toe, limiting running and sport
\end{itemize}

\section*{Prognosis and Outcomes}
Most broken toes heal well with simple care and leave no lasting problem. Pain settles within days to weeks, and the toe usually feels much better within four to six weeks, although mild swelling or stiffness may persist for several months. Big toe fractures take longer to settle because of the weight placed on them during walking. Smokers, older people, and those with diabetes or poor circulation can expect slower healing. Occasionally a toe heals slightly crooked, leaving a permanent bump or a subtle change in shape, which is generally harmless and does not affect function.

\section*{Prevention and Self-Care}
\begin{itemize}
    \item Wear protective footwear such as steel-capped boots when working with heavy objects
    \item Wear sturdy, closed-in shoes for sport and around the home, and avoid walking barefoot in busy or dim areas
    \item Keep floors clear of clutter and furniture that you are likely to strike
    \item Maintain good lighting at night to avoid stubbing toes
    \item Support bone health with adequate calcium and vitamin D and regular weight-bearing activity
    \item If you smoke, seek help to quit, since smoking slows bone healing
    \item After injury, follow the strapping, footwear, and follow-up advice given by your doctor
\end{itemize}

\section*{Sources and Further Reading}
The following organisations provide reliable information on toe fractures, foot injury first aid, and when to seek care.

\begin{itemize}
    \item NHS. \textit{Broken toe: symptoms, treatment, and self-care.} https://www.nhs.uk/conditions/broken-toe/
    \item Mayo Clinic. \textit{Broken toe: causes, diagnosis, and treatment.} https://www.mayoclinic.org/
    \item MedlinePlus. \textit{Toe injuries and fractures health information.} https://medlineplus.gov/
    \item MSD Manual. \textit{Fractures of the foot and toes.} https://www.msdmanuals.com/
    \item American Academy of Orthopaedic Surgeons. \textit{Toe and forefoot fractures patient education.} https://www.aaos.org/
    \item National Institute for Health and Care Excellence. \textit{Fracture assessment and management guidance.} https://www.nice.org.uk/
\end{itemize}
